% Asset ID: A21TrigonometricFunctionsTikZ-005
% Source: Chapter 6 reciprocal trig proof examples + CCEA A21-TRIG-LO008
% Related lesson section: Exam Technique and Worked Examples
% Purpose: Summarise the proof strategy used for reciprocal trig identities.

\documentclass[tikz,border=8pt]{standalone}
\usepackage{amsmath}
\usetikzlibrary{positioning, arrows.meta}
\begin{document}
\begin{tikzpicture}[step/.style={draw,rounded corners,align=center,minimum width=5.5cm,minimum height=1.05cm,inner sep=6pt},arrow/.style={-{Latex[length=3mm]},thick},warning/.style={draw,rounded corners,align=center,minimum width=6.7cm,minimum height=1.15cm,inner sep=6pt},title/.style={font=\large\bfseries}]
\node[title] (title) at (0,0) {Proof route for reciprocal trig identities};
\node[step, below=0.9cm of title] (one){Start with the messier side};
\node[step, below=0.65cm of one] (two){Rewrite reciprocal functions using sine and cosine\\[2pt]$\sec x=\dfrac{1}{\cos x}$,\quad $\cosec x=\dfrac{1}{\sin x}$,\quad $\cot x=\dfrac{\cos x}{\sin x}$};
\node[step, below=0.65cm of two] (three){Combine algebraic fractions using a common denominator};
\node[warning, below=0.65cm of three] (warn){\textbf{Warning:} cancel common factors only.\\Do not cancel across $+$ or $-$ signs.};
\node[step, below=0.65cm of warn] (four){Use identities such as\\[2pt]$\sin^2x+\cos^2x=1$,\\$1+\tan^2x=\sec^2x$,\\$1+\cot^2x=\cosec^2x$};
\node[step, below=0.65cm of four] (five){Simplify until the target side appears};
\node[step, below=0.65cm of five] (six){Conclude: LHS $\equiv$ RHS};
\draw[arrow] (one) -- (two);\draw[arrow] (two) -- (three);\draw[arrow] (three) -- (warn);\draw[arrow] (warn) -- (four);\draw[arrow] (four) -- (five);\draw[arrow] (five) -- (six);
\end{tikzpicture}
\end{document}
